\PassOptionsToPackage{table}{xcolor}
\documentclass[10pt,aspectratio=169]{beamer}
\newcommand{\LectureNumber}{13}
\input{course-theme.tex}
\title{String processing}
\subtitle{CSC103 Programming Fundamentals / Lecture 13}
\author{Dr. Muhammad Farhan}
\institute{Associate Professor of Computer Science\\Department of Computer Science\\COMSATS University Islamabad\\Sahiwal Campus}
\date{Fall 2026}
\begin{document}
\begin{frame}\titlepage\end{frame}

\begin{frame}[fragile,t]{The problem for this lecture}
\begin{columns}[T,onlytextwidth]\begin{column}{0.60\textwidth}
For a nonempty email-like string with one @, extract the username and domain. Test ali@campus.edu.
\end{column}\begin{column}{0.35\textwidth}\small
Find the @ position and reject missing or empty parts.
\end{column}\end{columns}
\note{Introduce the target problem before explaining the tools. Revisit it in the guided solution later.\par Syllabus reading: Deitel Ch 16. CLO-2.}
\end{frame}

\begin{frame}[fragile,t]{Learning objectives}
\begin{columns}[T,onlytextwidth]\begin{column}{0.60\textwidth}
\begin{itemize}
\item Use common String methods
\item Compare text content and extract substrings
\item Convert numeric text and numeric values
\end{itemize}
\end{column}\begin{column}{0.35\textwidth}\small
CLO-2

Length and character access\par Content comparison\par Substrings and searching\par Numbers and text
\end{column}\end{columns}
\note{Explain how each objective supports the opening problem. The final questions check these objectives.\par Syllabus reading: Deitel Ch 16. CLO-2.}
\end{frame}

\begin{frame}[fragile,t]{Length and character access}
\begin{itemize}
\item length() counts UTF-16 code units.
\item charAt(index) uses positions from 0 to length()-1.
\item Basic Latin examples let each char correspond to one visible character.
\end{itemize}
\vspace{1mm}\begin{block}{Example}
\begin{lstlisting}
String s = "Java";
s.length()    // 4
s.charAt(0)   // J
\end{lstlisting}
\end{block}
\note{charAt(length()) is outside the valid range. Supplementary Unicode characters may occupy two char positions.\par Syllabus reading: Deitel Ch 16. CLO-2.}
\end{frame}

\begin{frame}[fragile,t]{Content comparison}
\begin{itemize}
\item equals compares String content.
\item equalsIgnoreCase ignores case distinctions for its comparison.
\item compareTo returns a negative value, zero or a positive value.
\end{itemize}
\vspace{1mm}\begin{block}{Example}
\begin{lstlisting}
"cat".equals("cat")        // true
"Cat".equals("cat")        // false
"cat".compareTo("dog") < 0 // true
\end{lstlisting}
\end{block}
\note{== compares references, so string interning can make an incorrect approach appear to work. compareTo is lexicographic, not locale-aware dictionary sorting.\par Syllabus reading: Deitel Ch 16. CLO-2.}
\end{frame}

\begin{frame}[fragile,t]{Substrings and searching}
\begin{itemize}
\item substring(begin,end) includes begin and excludes end.
\item indexOf returns the first matching position or -1.
\item trim returns a result without leading or trailing basic whitespace.
\end{itemize}
\vspace{1mm}\begin{block}{Example}
\begin{lstlisting}
String code = "CS-103";
code.substring(3)     // "103"
code.substring(0, 2)  // "CS"
code.indexOf("-")     // 2
\end{lstlisting}
\end{block}
\note{State the valid condition 0 \textless{}= begin \textless{}= end \textless{}= length. Assign the returned result if the program needs to retain a transformed String.\par Syllabus reading: Deitel Ch 16. CLO-2.}
\end{frame}

\begin{frame}[fragile,t]{Numbers and text}
\begin{itemize}
\item Integer.parseInt and Double.parseDouble parse numeric text.
\item String.valueOf converts a value to text.
\item Malformed numeric text raises NumberFormatException.
\end{itemize}
\vspace{1mm}\begin{block}{Example}
\begin{lstlisting}
int n = Integer.parseInt("42");
String label = String.valueOf(n);
\end{lstlisting}
\end{block}
\note{Parsing "42.5" as int fails. Formatting is a separate concern from conversion. Exception handling is taught in Lectures 24 and 25.\par Syllabus reading: Deitel Ch 16. CLO-2.}
\end{frame}

\begin{frame}[fragile,t]{Substring boundaries form a half-open interval}
\begin{itemize}
\item A substring starts at begin and stops before end.
\item The difference end - begin gives its UTF-16 length.
\item An empty substring is valid when begin equals end.
\end{itemize}
\vspace{1mm}\begin{block}{Example}
\begin{lstlisting}
"JAVA".substring(1,3) is "AV"
Included indexes: 1 and 2
Excluded index: 3
\end{lstlisting}
\end{block}
\note{Explain the mechanism using the displayed example. A substring starts at begin and stops before end. The difference end - begin gives its UTF-16 length. An empty substring is valid when begin equals end.\par Syllabus reading: Deitel Ch 16. CLO-2.}
\end{frame}

\begin{frame}[fragile,t]{Search failure must be handled before slicing}
\begin{itemize}
\item indexOf returns -1 when it finds no match.
\item Passing -1 as a substring boundary is invalid.
\item Validate the position before extracting the two parts of a structured value.
\end{itemize}
\vspace{1mm}\begin{block}{Example}
\begin{lstlisting}
For "ali@campus.edu": position=3
Username: substring(0,3)
Domain: substring(4)
\end{lstlisting}
\end{block}
\note{Explain the mechanism using the displayed example. indexOf returns -1 when it finds no match. Passing -1 as a substring boundary is invalid. Validate the position before extracting the two parts of a structured value.\par Syllabus reading: Deitel Ch 16. CLO-2.}
\end{frame}


\begin{frame}[fragile,t]{Substring boundaries sit between characters}
\begin{center}\begin{tikzpicture}
\node[box,text width=13mm,minimum height=11mm] (n0) at (0.0,0.0) {C};
\node[font=\scriptsize,text=accent] at (0.0,.85) {index 0};
\node[box,text width=13mm,minimum height=11mm] (n1) at (1.8,0.0) {S};
\node[font=\scriptsize,text=accent] at (1.8,.85) {index 1};
\node[box,text width=13mm,minimum height=11mm] (n2) at (3.6,0.0) {C};
\node[font=\scriptsize,text=accent] at (3.6,.85) {index 2};
\node[box,text width=13mm,minimum height=11mm] (n3) at (5.4,0.0) {1};
\node[font=\scriptsize,text=accent] at (5.4,.85) {index 3};
\node[box,text width=13mm,minimum height=11mm] (n4) at (7.2,0.0) {0};
\node[font=\scriptsize,text=accent] at (7.2,.85) {index 4};
\node[box,text width=13mm,minimum height=11mm] (n5) at (9.0,0.0) {3};
\node[font=\scriptsize,text=accent] at (9.0,.85) {index 5};
\draw[very thick,draw=berry](-.75,-.8)--(4.35,-.8);\node[text=berry,font=\small] at (1.8,-1.25){substring(0,3)};
\end{tikzpicture}\end{center}
\begin{block}{What to notice}substring(0,3) selects positions 0, 1 and 2. The end index 3 is excluded.\end{block}
\note{Trace each labelled connection. substring(0,3) selects positions 0, 1 and 2. The end index 3 is excluded.}
\end{frame}
\begin{frame}[fragile,t]{Key distinctions}
\small\renewcommand{\arraystretch}{1.5}
\rowcolors{2}{pale}{white}
\begin{tabularx}{\textwidth}{>{\raggedright\arraybackslash}X>{\raggedright\arraybackslash}X>{\raggedright\arraybackslash}X}
\rowcolor{navy}
\color{white}\bfseries Operation & \color{white}\bfseries Question answered & \color{white}\bfseries Example result\\
equals & Same text content? & "CS" equals "CS": true\\
== & Same object reference? & May differ for equal text\\
indexOf & Where is the first match? & "CS-103" finds "-" at 2\\
substring & Which interval of text? & substring(3) gives "103"\\
\end{tabularx}
\vspace{3mm}\footnotesize These distinctions determine which operation or control structure fits the problem.
\note{\par Syllabus reading: Deitel Ch 16. CLO-2.}
\end{frame}

\begin{frame}[fragile,t]{First example: execution}
\begin{lstlisting}
String code = "CS-103";
String prefix = code.substring(0, 2);
int number = Integer.parseInt(code.substring(3));
System.out.println(prefix);
System.out.println(number + 1);
\end{lstlisting}
\note{This is the main body from Java\_Examples/Lecture13.java. The source file includes the complete class. Predict the result before the next slide.\par Syllabus reading: Deitel Ch 16. CLO-2.}
\end{frame}

\begin{frame}[fragile,t]{First example: result}
\begin{columns}[T,onlytextwidth]\begin{column}{0.60\textwidth}
CS\par 104\par The numeric substring is "103". Parsing enables arithmetic.
\end{column}\begin{column}{0.35\textwidth}\small
Substrings and searching

Numbers and text
\end{column}\end{columns}
\note{Expected output and interpretation: CS
104
The numeric substring is "103". Parsing enables arithmetic.\par Syllabus reading: Deitel Ch 16. CLO-2.}
\end{frame}

\begin{frame}[fragile,t]{Guided practice}
\begin{columns}[T,onlytextwidth]\begin{column}{0.60\textwidth}
For a nonempty email-like string with one @, extract the username and domain. Test ali@campus.edu.
\end{column}\begin{column}{0.35\textwidth}\small
Attempt the solution before continuing.

Use the definitions and examples from this lecture.
\end{column}\end{columns}
\note{Give students 8 minutes to attempt the problem. The following slides provide a complete solution. Original solution guidance: Find position with indexOf("@"). Use substring(0,pos) and substring(pos+1). Results: ali and campus.edu. Reject a missing @ before extracting.\par Syllabus reading: Deitel Ch 16. CLO-2.}
\end{frame}

\begin{frame}[fragile,t]{Solution strategy}
\begin{columns}[T,onlytextwidth]\begin{column}{0.60\textwidth}
\begin{itemize}
\item Find the @ position and reject missing or empty parts.
\item Extract text before @ as the username.
\item Extract text after @ as the domain.
\end{itemize}
\end{column}\begin{column}{0.35\textwidth}\small
The next program uses fixed inputs so its result can be reproduced.
\end{column}\end{columns}
\note{Discuss why each step is needed. State the input assumptions before executing the program.\par Syllabus reading: Deitel Ch 16. CLO-2.}
\end{frame}

\begin{frame}[fragile,t]{Complete solution}
\begin{lstlisting}
public class Solution13 {
  public static void main(String[] args) throws Exception {
    String email = "ali@campus.edu";
    int at = email.indexOf("@");
    if (at <= 0 || at == email.length() - 1) {
      System.out.println("Invalid format");
    } else {
      String user = email.substring(0, at);
      String domain = email.substring(at + 1);
      System.out.println(user);
      System.out.println(domain);
    }
  }
}
\end{lstlisting}
\note{Complete runnable file: Java\_Examples/Solution13.java. Inputs are fixed in the program.  \par Syllabus reading: Deitel Ch 16. CLO-2.}
\end{frame}

\begin{frame}[fragile,t]{Execution trace}
\small\renewcommand{\arraystretch}{1.5}
\rowcolors{2}{pale}{white}
\begin{tabularx}{\textwidth}{>{\raggedright\arraybackslash}X>{\raggedright\arraybackslash}X>{\raggedright\arraybackslash}X}
\rowcolor{navy}
\color{white}\bfseries Part & \color{white}\bfseries Index interval & \color{white}\bfseries Result\\
Username & [0, 3) & ali\\
Separator & index 3 & @\\
Domain & [4, 14) & campus.edu\\
\end{tabularx}
\vspace{3mm}\footnotesize Follow the changing state in execution order.
\note{Manually derived trace for the complete solution. Find the @ position and reject missing or empty parts. Extract text before @ as the username. Extract text after @ as the domain.\par Syllabus reading: Deitel Ch 16. CLO-2.}
\end{frame}

\begin{frame}[fragile,t]{Solution result}
\begin{columns}[T,onlytextwidth]\begin{column}{0.60\textwidth}
ali\par campus.edu
\end{column}\begin{column}{0.35\textwidth}\small
Why this result follows

Extract text after @ as the domain.
\end{column}\end{columns}
\note{Expected result: ali
campus.edu\par Syllabus reading: Deitel Ch 16. CLO-2.}
\end{frame}

\begin{frame}[fragile,t]{A common incorrect approach}
String name = new String("Ali");\par if (name == "Ali") \{ /* matched */ \}
\vspace{1mm}\begin{block}{Example}
\begin{lstlisting}
Use name.equals("Ali") or "Ali".equals(name). == checks reference identity, not text content.
\end{lstlisting}
\end{block}
\note{Ask students to predict the failure before discussing the explanation. The right side gives the correction.\par Syllabus reading: Deitel Ch 16. CLO-2.}
\end{frame}

\begin{frame}[fragile,t]{Boundary and failure checks}
\begin{columns}[T,onlytextwidth]\begin{column}{0.60\textwidth}
Check the stated input assumptions.

Read an item code shaped like BOOK-250. Extract the category and price, then print the price plus 10.
\end{column}\begin{column}{0.35\textwidth}\small
Find position with indexOf("@"). Use substring(0,pos) and substring(pos+1). Results: ali and campus.edu. Reject a missing @ before extracting.
\end{column}\end{columns}
\note{Use the specification to derive each expected result. Exercise solution: Find position with indexOf("@"). Use substring(0,pos) and substring(pos+1). Results: ali and campus.edu. Reject a missing @ before extracting.\par Syllabus reading: Deitel Ch 16. CLO-2.}
\end{frame}

\begin{frame}[fragile,t]{Check your understanding}
\begin{columns}[T,onlytextwidth]\begin{column}{0.60\textwidth}
What is "JAVA".substring(1,3)? Does s.trim() change the String referred to by s?
\end{column}\begin{column}{0.35\textwidth}\small
Write a prediction and explain the rule that supports it.
\end{column}\end{columns}
\note{Answer: "AV". No. Use s = s.trim() to keep the returned result.\par Syllabus reading: Deitel Ch 16. CLO-2.}
\end{frame}

\begin{frame}[fragile,t]{Answer and explanation}
\begin{columns}[T,onlytextwidth]\begin{column}{0.60\textwidth}
"AV". No. Use s = s.trim() to keep the returned result.
\end{column}\begin{column}{0.35\textwidth}\small
Use name.equals("Ali") or "Ali".equals(name). == checks reference identity, not text content.
\end{column}\end{columns}
\note{Reconnect the answer to the relevant definition rather than treating it as a fact to memorise.\par Syllabus reading: Deitel Ch 16. CLO-2.}
\end{frame}


\begin{frame}[fragile,t]{Compare cases and predict outcomes}
\small\renewcommand{\arraystretch}{1.5}\rowcolors{2}{pale}{white}
\begin{tabularx}{\textwidth}{>{\raggedright\arraybackslash}X>{\raggedright\arraybackslash}X>{\raggedright\arraybackslash}X}\rowcolor{navy}
\color{white}\bfseries Call on "CSC103" & \color{white}\bfseries Indexes used & \color{white}\bfseries Result\\
substring(0,3) & 0, 1, 2 & CSC\\
substring(3) & 3, 4, 5 & 103\\
substring(6) & No characters & Empty string\\
\end{tabularx}\par\vspace{3mm}\begin{exampleblock}{Discuss}Explain each expected result before running the example. Which case would reveal a wrong assumption?\end{exampleblock}
\note{Read the first two columns and invite predictions before discussing the outcome.}
\end{frame}

\begin{frame}[fragile,t]{Apply the idea}
\begin{alertblock}{Independent practice}Given "lab:42", find the colon and convert the text after it to an int. State the assumptions needed for the conversion.\end{alertblock}\vspace{3mm}\begin{enumerate}\item Work through the input by hand.\item Write or correct the program.\item Justify the expected result.\end{enumerate}\note{Allow 3–5 minutes, then compare answers in pairs. Given "lab:42", find the colon and convert the text after it to an int. State the assumptions needed for the conversion.}
\end{frame}

\begin{frame}[fragile,t]{Solution and reasoning}
\begin{lstlisting}
String text = "lab:42";
int separator = text.indexOf(":");
String number = text.substring(separator + 1);
int value = Integer.parseInt(number);
System.out.println(value + 1);
\end{lstlisting}\vspace{1mm}\small\begin{itemize}
\item The colon is at index 3; substring(4) is "42".
\item Parsing gives 42, so the output is 43.
\item This example assumes a colon and a valid integer suffix; production code must validate them.
\end{itemize}\note{Answer: The colon is at index 3; substring(4) is "42". Parsing gives 42, so the output is 43. This example assumes a colon and a valid integer suffix; production code must validate them.}
\end{frame}

\begin{frame}[fragile,t]{Counting vowels in a sentence}
\begin{itemize}
\item Scan each character once and increase the count only on a vowel match.
\item Normalise case to make uppercase and lowercase English vowels equivalent.
\item This exercise counts a, e, i, o and u; it is not a general linguistic vowel detector.
\end{itemize}
\vspace{3mm}\footnotesize\textcolor{accent}{Source: supplied ISB campus slides. See speaker notes and ISB\_Source\_Map.md.}
\note{Adapted from the supplied ISB campus folder: Strings {-}practice.pptx, slides 8, 9; Methods Practice.pptx, slides 10, 11. Uses Locale.ROOT and a compact membership check while retaining the source phrase and vowel{-}count task. Source exercise deck credits Instructor Khurram Iqbal.}
\end{frame}

\begin{frame}[fragile,t]{Worked problem: Counting vowels in a sentence}
\begin{block}{Task}Count the vowels in Welcome to Java. Predict the six matching characters before running.\end{block}
\small Input: Values are fixed in the program for a reproducible demonstration.\par\vspace{3mm}\begin{exampleblock}{Before running}Predict the output and identify the variables that change.\end{exampleblock}
\note{Adapted from the supplied ISB campus folder: Strings {-}practice.pptx, slides 8, 9; Methods Practice.pptx, slides 10, 11. Uses Locale.ROOT and a compact membership check while retaining the source phrase and vowel{-}count task. Source exercise deck credits Instructor Khurram Iqbal.}
\end{frame}

\begin{frame}[fragile,t]{Complete ISB example (1/1)}
\begin{lstlisting}
public class IsbLecture13 {
  public static void main(String[] args) throws Exception {
    String text = "Welcome to Java".toLowerCase(java.util.Locale.ROOT);
    int count = 0;
    for (int i = 0; i < text.length(); i++) {
      char ch = text.charAt(i);
      if ("aeiou".indexOf(ch) >= 0) count++;
    }
    System.out.println(count);
  }

}
\end{lstlisting}
\note{Adapted from the supplied ISB campus folder: Strings {-}practice.pptx, slides 8, 9; Methods Practice.pptx, slides 10, 11. Uses Locale.ROOT and a compact membership check while retaining the source phrase and vowel{-}count task. Source exercise deck credits Instructor Khurram Iqbal. Complete file: ISB\_Java\_Examples/IsbLecture13.java.}
\end{frame}

\begin{frame}[fragile,t]{Worked trace and expected result}
\small\renewcommand{\arraystretch}{1.4}\rowcolors{2}{pale}{white}
\begin{tabularx}{\textwidth}{>{\raggedright\arraybackslash}X>{\raggedright\arraybackslash}X>{\raggedright\arraybackslash}X}\rowcolor{navy}
\color{white}\bfseries Word & \color{white}\bfseries Matching vowels & \color{white}\bfseries Cumulative count\\
Welcome & e, o, e & 3\\
to & o & 4\\
Java & a, a & 6\\
\end{tabularx}\par\vspace{2mm}
\begin{block}{Expected console output}\ttfamily\footnotesize 6\end{block}
\note{Adapted from the supplied ISB campus folder: Strings {-}practice.pptx, slides 8, 9; Methods Practice.pptx, slides 10, 11. Uses Locale.ROOT and a compact membership check while retaining the source phrase and vowel{-}count task. Source exercise deck credits Instructor Khurram Iqbal. Trace and expected values independently worked for this edition.}
\end{frame}

\begin{frame}[fragile,t]{Extend the ISB example}
\begin{alertblock}{Practice}Predict the counts for an empty string, AEIOU and rhythm.\end{alertblock}\vspace{3mm}\begin{exampleblock}{Explain your reasoning}0, 5 and 0 under the stated five{-}vowel contract.\end{exampleblock}
\note{Adapted from the supplied ISB campus folder: Strings {-}practice.pptx, slides 8, 9; Methods Practice.pptx, slides 10, 11. Uses Locale.ROOT and a compact membership check while retaining the source phrase and vowel{-}count task. Source exercise deck credits Instructor Khurram Iqbal. Answer: 0, 5 and 0 under the stated five{-}vowel contract.}
\end{frame}
\begin{frame}[fragile,t]{Lecture summary}
\begin{columns}[T,onlytextwidth]\begin{column}{0.60\textwidth}
\begin{itemize}
\item common String methods
\item Compare text content and extract substrings
\item Convert numeric text and numeric values
\end{itemize}
\end{column}\begin{column}{0.35\textwidth}\small
\begin{itemize}
\item Find the @ position and reject missing or empty parts.
\item Extract text before @ as the username.
\item Extract text after @ as the domain.
\end{itemize}
\end{column}\end{columns}
\note{Ask students to give one concrete example for the concepts listed. Use the worked solution to connect the topics.\par Syllabus reading: Deitel Ch 16. CLO-2.}
\end{frame}

\begin{frame}[fragile,t]{After-class practice}
\begin{columns}[T,onlytextwidth]\begin{column}{0.60\textwidth}
Read an item code shaped like BOOK-250. Extract the category and price, then print the price plus 10.
\end{column}\begin{column}{0.35\textwidth}\small
Syllabus reading\par Deitel Ch 16

Next lecture: Built-in and instance methods
\end{column}\end{columns}
\note{Reading labels follow the supplied outline. Check topic headings against the available textbook edition. Require a program and expected results for the practice task.\par Syllabus reading: Deitel Ch 16. CLO-2.}
\end{frame}
\end{document}
